\section{The exact cohomology sequence}
\label{XI.4}

Let $S$ be a prescheme, so that the category
$\Sch_{/S}$ of preschemes over~$X$ is determined,
hence also the notion of group in it, which will also be called a
\emph{group prescheme over~$S$}, or simply an
\emph{$S$-group}. To simplify the exposition and fix
ideas, we shall most often restrict ourselves in what follows to
groups that are \emph{affine} and \emph{flat} over~$S$\kern1pt\footnote{In
fact, for what follows, the quasi-affine hypothesis instead
of the affine one would suffice, \cf footnote~\eqref{note-296}, p\ptbl\pageref{note-296}.}, which will
suffice for the applications we have in view. (Of course, one
encounters numerous cases where neither the one nor the other hypothesis
is satisfied). Let
% original p. 293
$G$ be such an $S$-group, and let $P$ be a prescheme over~$S$ on
which $G$ acts on the right, which implies in particular a
morphism
$$
\pi:P\times_SG\to P
$$
in the category $\Sch_{/S}$, satisfying the well-known
axioms. One says that $P$ is \emph{formally principal homogeneous
under} $G$ if the morphism
$$
P\times_S G\to P\times_S P
$$
with components $\pr_1$ and $\pi$ is an isomorphism; it amounts to the
same to say that for every object $S'$ of~$\Sch_{/S'}$,
$P(S')=\Hom_S(S',P)$, considered as a set with groups
of operators $G(S')=\Hom_S(S',G)$, is empty or principal
homogeneous (\ie empty or isomorphic to $G(S')$ on which the
group $G(S')$ acts by right translations). One says that
$P$ is \emph{trivial} if $P$ is isomorphic to $G$, on which $G$
acts by right translations, or what amounts to the
same, if each of the sets with operators $P(S')$ under
$G(S')$ is trivial. One verifies, for example by the
patented procedure of passing to the set-theoretic case, that $P$
\emph{is trivial if and only if it
is formally principal homogeneous, and admits a section over~$S$}
(this last fact being stated in categorical terms by
saying that $P$ has a section over the final object $e=S$ of~$\Sch_{/S}$,
\ie that there exists a morphism from~$e$ to $P$). To define the
notion of principal homogeneous bundle $P$ under $G$, stronger
than that of formally principal homogeneous bundle, one must
first specify in $\Sch_{/S}$ a set of morphisms that
will be used for ``descent'', and will play the role of
``localization morphisms'' in order to ``trivialize'' bundles. The
most suitable choice varies with the context, none
containing all the others\footnote{\label{note-293}See on this subject
SGA~3~IV, in particular~\S 4.}. Here it will be convenient to adopt the
following definition:

\begin{definition}
\label{XI.4.1}
Let $G$ be an $S$-group. One calls a \emph{principal
homogeneous bundle} (on the right) under~$G$ an $S$-prescheme $P$
with $S$-group~$G$ acting on the right, such that there exist a covering of~$S$ by open sets~$U_i$, and for every $i$ a faithfully flat
and quasi-compact change of base morphism $S'_i\to U_i$, such that
$P'=P\times_S S'$ is a trivial prescheme with operators
under $G'=G\times_S S'$, where $S'$ is the
$S$-prescheme disjoint sum of the~$S'_i$.
\end{definition}

(Note
% original p. 294
that the change of base functor $X\mto X'=X\times_S
S'$,
being left exact, transforms groups into groups, objects
with group of operators into objects with groups
of operators). Let us note that \Ref{XI.4.1} is \emph{stable under
change of base}. Let us also note:

\begin{proposition}
\label{XI.4.2}
Let $G$ be an $S$-group, flat and quasi-compact over~$S$, and $P$ an
$S$-prescheme on which $G$ acts on the right. The following conditions
are equivalent:
\begin{enumerate}
\item[(i)] $P$ is a principal homogeneous bundle
under~$G$;
\item[(ii)] $P$ is formally principal homogeneous under~$G$,
and the structure morphism $P\to S$ is faithfully flat and
quasi-compact.
\end{enumerate}
\end{proposition}

If $P$ is principal homogeneous under~$G$, then with the notation
of~\Ref{XI.4.1}, $P'$ is faithfully flat and quasi-compact over~$S'$
(since $G'$ is, and $P'$ is $S'$-isomorphic to it), hence $P$ has the
same properties over~$S$ (for ``surjective'' and
``quasi-compact'', \cf VIII~\Ref{VIII.3.1}; for ``flat'' it is an
omission in the sorites of Expos\'e~\Ref{VIII}). Conversely, if
(ii) is satisfied, let us take the change of base $S'=P$, which
is indeed faithfully flat and quasi-compact over~$S$; then $P'$ will be
formally principal homogeneous over~$S'$ since $P$ is so over~$S$
and the change of base functor is left exact; on the other
hand $P'$ has a section over~$S'$,
namely the diagonal section, hence
it is a trivial principal bundle, which completes the
proof.

\begin{corollary}
\label{XI.4.3}
If $G$ is affine and flat over~$S$, every principal
homogeneous bundle $P$ under $G$ is affine and flat over~$S$.
\end{corollary}

Indeed, it becomes so after a faithfully flat and
quasi-compact extension of the base, and one applies (VIII~\Ref{VIII.5.6}).

The usefulness of Definition~\Ref{XI.4.1} for $S$-groups that are
\emph{flat} and \emph{affine} over~$S$ is due to
(VIII~\Ref{VIII.2.1}), \ie to the fact that the morphisms $S'\to S$
considered in~\Ref{XI.4.1} are effective descent morphisms
for the category of preschemes affine over
others. Thanks to this fact, the verification of the facts
sketched below is done in an essentially
``categorical'' way\footnote{\Cf \loccit in
footnote~\eqref{note-293}, p\ptbl\pageref{note-293}.}. Let $E$ be an $S$-prescheme on which the
$S$-group $G$ acts on the left, and let $P$ be a principal
homogeneous bundle (on the right) under~$G$; we want to
define an associated bundle $E^{(P)}$, ``locally''
isomorphic to~$E$. For this, let us make $G$ act
% original p. 295
on the right on $P\times_S E$ according to the law $(x,y)\mto
(xg,g^{-1}y)$, which describes such operations in the
set-theoretic context, and extends to categories by the
patented procedure. We shall set, subject to existence:
$$
E^{(P)}=(P\times_S E)/G
$$
whereupon one finds that $P\times_S E$ will be a prescheme
over~$T=E^{(P)}$, with right groups of operators
$G_T=G\times_S T$; to be at ease, one would like moreover
$P\times_S E$ to be a principal homogeneous bundle over~$T$ with
group $G_T$. To verify the existence of~$E^{(P)}$ and the
preceding property, let us take again the $S'$ of
Definition~\Ref{XI.4.1} and look at the situation over~$S'$
that is the inverse image of the initial situation. From the fact that $P'$ is
trivial,
\ie isomorphic to $G'_d$, one sees at once that $E^{\prime(P')}$
exists, and has the desired exactness property. In fact,
$E'\times_{S'} P'$ is $G'$-isomorphic to the product $E'\times_{S'} G'$,
hence $E^{\prime(P')}$ is isomorphic to~$E'$. Moreover, the formation of the
``associated bundle'' in the case of a trivial space with operators
commutes with every extension of the base, and taking in this
instance the various extensions of the base
$\xymatrix@C=.5cm{S''\ar@<2pt>[r]\ar@<-2pt>[r]&S'}$
and
$\xymatrix@C=.5cm{S'''\ar@<4pt>[r]\ar@<0pt>[r]\ar@<-4pt>[r]&S'}$, where
$S''$ and $S'''$ are the double and triple fiber products of~$S'$
over~$S$, one finds that $E^{\prime(P')}$ \emph{is endowed with a descent
datum relative to the morphism $S'\to S$, and that $E^{(P)}$
exists with the required properties if and only if this
descent datum is effective}; of course $E^{(P)}$ is then
nothing but the descended object. (Use the fact that $S'\to S$ is
a descent morphism in the category of
$S$-preschemes, \cf VIII~\Ref{VIII.5.2}). It follows~that
\emph{the associated bundle exists if $E$ is affine over~$S$}. We
shall apply this construction to the case where one has a homomorphism of~$S$-groups $G\to H$, and where one takes for $E$ the $S$-prescheme~$H$ endowed with the left operations of~$G$ on~$H$ resulting
from the given morphism; since $H$~acts on the right on
itself in such a way as to commute with the operations of~$G$
on~$H$, and since (subject to existence over~$S$) the
formation of the associated bundle commutes with extension of the
base, one easily finds that $H$ will act on the right on~$P^{(H)}$, which is therefore a principal homogeneous bundle
under $H$ in the sense of~\Ref{XI.4.1}, and
% original p. 296
more precisely is trivialized by the same
morphism $S'\to S$ as~$P$. In particular, \emph{to every
principal homogeneous bundle $P$ under $G$ and every homomorphism of~$S$-groups $G\to H$, with $H$ affine over~$S$, is associated a
principal homogeneous bundle with group~$H$}, in a way
functorial in $(G\to H)$, and compatible with arbitrary changes of base
$T\to S$.

\begin{definition}
\label{XI.4.4}
Let $G$ be an $S$-prescheme. One denotes by $\H^0(S,G)$ the set of
sections of~$G$ over~$S$, which one will regard as a group when
$G$ is an $S$-group. In this case, one denotes by $\H^1(S,G)$ the set of
classes, up to isomorphism, of principal homogeneous
bundles over~$S$ with group $S$, regarding $\H^1(S,G)$
as endowed with the ``marked point'' that corresponds to the trivial
bundles\footnote{\label{note-296}This notation is consistent
with the general cohomological notation (SGA~4~V)
only when one has criteria for effectiveness of descent,
which are hardly ensured except if $G$ is affine (or only
quasi-affine, \cf (VIII~\Ref{VIII.7.9})).}.
\end{definition}

Thus, $\H^0(S,G)$ is a functor in the $S$-prescheme~$G$,
with values in the category of sets. This functor is
left exact, a fortiori commutes with finite products, which
indeed implies that it transforms groups into groups, commutative
groups into commutative groups. Similarly,
$\H^1(S,G)$ is a functor in the \emph{affine} $S$-group~$G$, with
values in the category of sets, thanks to the
formation of associated bundles; one easily finds that this
functor commutes with finite products. In particular it transforms the
groups in the category of affine $S$-groups, \ie the
\emph{commutative affine} $S$-groups, into groups of the second,
and even into commutative groups (since the groups of the
first category are commutative). Thus, \emph{if $G$ is
a commutative affine $S$-group, $\H^1(S,G)$ is a commutative
group}, and a homomorphism $G\to H$ of commutative affine $S$-groups
gives rise to a homomorphism of groups
$\H^1(S,G)\to\H^1(S,H)$.

For simplicity, we restrict ourselves in what follows to the
consideration of \emph{affine and commutative} $S$-groups. Let
$$
0\to G'\lto{u} G\lto{v}G''\to 0
$$
be a sequence of morphisms of such groups; \emph{we shall say that this
sequence is exact if $vu=0$} (which allows one to regard $G$
as a prescheme over~$G''$, with right groups of operators~$G'_{G''}$) \emph{and if $G$ is a principal homogeneous
bundle
% original p. 297
over~$G''$ with group $G'_{G''}=G'\times_S G''$}. This implies in
particular that $u\colon G'\to G$ is a kernel of~$v$, and a fortiori
it implies the exactness of the sequence
$$0\to\H^0(X,G')\to\H^0(X,G)\to\H^0(X,G'').$$ It implies moreover the
possibility of defining a map
$$
\partial\colon\H^0(X,G'')\to\H^1(X,G')\quoi,
$$
by associating to every section of~$G''$ over~$S$, \ie to every
$S$-morphism $f\colon S\to G''$, the principal homogeneous bundle
$P_f$ with group $G'\simeq f^*(G'_{G''})$ over~$S$, inverse image of the
principal homogeneous bundle $G$ over~$G''$. From the point of view of
$S$-preschemes, it is therefore nothing but the inverse image under
$v\colon G\to G''$ of the subprescheme image of~$S$ under
the immersion~$f$, and the operations of~$G'$ on~$P_f$ are induced
by the right operations of~$G'$
on~$G$.
We likewise leave to the reader the verification of the
following proposition, which presents no difficulties
other than those of writing it up:

\begin{proposition}
\label{XI.4.5}
The map $\partial\colon \H^0(X,G'')\to\H^1(X,G')$ is a
homomorphism of groups. The following sequence of homomorphisms is
exact:
\begin{multline*}
0\to\H^0(X,G')\lto{u^0}\H^0(X,G)\lto{v^0}\H^0(X,G')\lto{\partial}\H^1(X,G'')\\
\lto{u^1}\H^1(X,G) \lto{v^1}\H^1(X,G'')
\end{multline*}
(where the homomorphisms other than $\partial$ come from the
functoriality of~$\H^0$ \resp~$\H^1$).
\end{proposition}

\begin{remarks}
\label{XI.4.6}
The point of view presented here for the study of principal
homogeneous bundles is visibly inspired by Serre~\cite{XI.7}, which
the reader will have every interest in consulting. When one wants
a formalism that applies also to structural $S$-groups
that are quasi-projective over~$S$ (so as to
include projective abelian schemes in particular), it is
advantageous to modify \Ref{XI.4.1} by requiring there that $S'$
be a sum of preschemes $S'_i$ that are finite and locally
free over open sets $S_i$ of~$S$ covering $S$. The
preceding developments are then valid, including
in particular \Ref{XI.4.5}, on replacing everywhere the affine hypothesis
by the quasi-projective hypothesis, and on interpreting
correspondingly the definition given above of an
exact sequence of~$S$-groups. It suffices indeed to replace the
reference to (VIII~\Ref{VIII.2.1})
% original p. 298
by (VIII~\Ref{VIII.7.7}): the morphisms $S'\to S$ used are
effective descent morphisms for the fibered category
of preschemes quasi-projective over others. One should
take care, however, that this second notion of principal
homogeneous bundle is more restrictive than the
first~\Ref{XI.4.1}.
\end{remarks}

\subsection{}
\label{XI.4.7}
One obtains an even more restrictive notion of principal
homogeneous bundle by requiring that $S$ be covered by open sets $S_i$
such that for every~$i$, $P|S_i$ is a trivial bundle with operators
under~$G|S_i$: one will then say that $P$ is a \emph{locally trivial}
principal homogeneous bundle.
The classes of these bundles, for given $G$, form a subset
of~$\H^1(X,G)$, which is in one-to-one correspondence with
$\H^1(X,\cal{O}_X(G))$, where $\cal{O}_X(G)$ is the sheaf (in the
ordinary sense) of sections of~$G$ over~$S$, \cf \cite{XI.2}. For these $\H^1$
to still give rise to an exact cohomology
sequence~\Ref{XI.4.5}, one must obviously assume that the sequence
$0\to G'\to G\to G''\to 0$ is exact in the reasonable sense for this
new context, \ie that $G$ is a locally trivial bundle
over~$G''$, with group $G'_{G''}$; this also means that $u\colon
G'\to G$ is a kernel of~$v\colon G\to G''$, and that $G$ locally admits
a section over~$G''$.

\subsection{}
\label{XI.4.8}
It is obviously very desirable to continue the exact
sequence~\Ref{XI.4.5} by introducing the higher cohomology groups
$\H^i(X,G)$. This is possible by placing oneself in
the framework of ``Weil cohomology'': one considers the
category $\cal{B}$ of quasi-compact preschemes over~$S$,
endowed with the set $\cal{M}$ of faithfully flat and
quasi-compact morphisms, which will be called localizing morphisms. An abelian ``Weil
sheaf'' on~$S$ (or better, on~$(\cal{B},\cal{M})$) is
then a contravariant functor $\cal{F}$ from~$\cal{B}$ to the
category of abelian groups, transforming sums into
products, and a sequence
${\def\labelstyle{\scriptstyle}\xymatrix@C=.5cm{T''=T'\times_TT'\ar@<+2pt>[rr]^-{\pr_1,\pr_2}\ar@<-2pt>[rr] &&T'\ar[r]^f & T}}$, with $f\in\cal{M}$, into an \emph{exact} diagram
of sets $\xymatrix@C=.5cm{\cal{F}(T)\ar[r] & \cal{F}(T')
\ar@<2pt>[r]\ar@<-2pt>[r] & \cal{F}(T'')}$. The Weil sheaves
form an abelian category with exact inductive limits
admitting a generator, hence admitting enough
injective objects~\cite{XI.1}. The right derived functors of the
functor $\Gamma(\cal{F})=\cal{F}(S)$ are then denoted
$\H^i(S,\cal{F})$. On the other hand, every commutative $S$-group
obviously defines a Weil sheaf (VIII~\Ref{VIII.5.2}),
whose $\H^0$ and $\H^1$ are none other than $\H^0(S,G)$ and
$\H^1(S,G)$, which allows one to define the other $\H^i(S,G)$ in a
reasonable way. One
% original p. 299
shows moreover that an exact sequence of~$S$-groups defines an
exact sequence of Weil sheaves, which allows one to recover and to
extend the exact sequence~\Ref{XI.4.5}\footnote{For a
systematic study of this point of view, \cf SGA~4~I~to~IX.}.

\subsection{}
\label{XI.4.9}
It would be advisable to develop the noncommutative variants
of~\Ref{XI.4.5} as in~\cite{XI.2}. For a systematic
development, in the appropriate framework, of the various
cohomological notions sketched in the present no., we
refer to a work in preparation by
J\ptbl \textsc{Giraud}\footnote{\Cf J\ptbl \textsc{Giraud},
\emph{Cohomologie non ab\'elienne},
Springer-Verlag 1971.}.

\section{Special cases of principal bundles}
\label{XI.5}

Suppose now that $S$ is connected and endowed with a geometric
point~$a$, whence a fundamental group $\pi_1(S,a)$
making it possible to classify the \'etale coverings of~$S$: the
category of \'etale coverings of~$S$ is
equivalent to the category of finite sets on which
$\pi_1$ acts continuously. It follows that a finite and \'etale group
scheme $G$ over~$S$ is
essentially determined by an ordinary finite group $\cal{G}$, on which
$\pi_1$ acts continuously by group automorphisms. An
\'etale covering $P$ of~$S$ on which $G$ acts on the right
is essentially determined by a finite set $\cal{P}$
on which $\pi_1$ acts continuously (on the left), and on which
$\cal{G}$ acts on the right in a way compatible with the
operations of~$\pi_1$:
$$
s(p\cdot g)=(sp)\cdot(sg)\qquad \text{for }s\in\pi_1,\
p\in\cal{P},\ g\in\cal{G}.
$$
One verifies that $P$ is a principal homogeneous bundle in the sense
of~\Ref{XI.4.1} if and only if $\cal{P}$ is a principal
homogeneous set under~$\cal{G}$ (use for example
criterion~\Ref{XI.4.2}). In other words, \emph{the category
of principal homogeneous bundles over~$S$ with group $G$ is
equivalent to the category of principal
homogeneous bundles with group $G$ in the category of finite sets
on which $\pi_1$ acts continuously.} One deduces from this in
particular a canonical bijection, functorial in $G$:
\begin{equation*}
\label{eq:XI.5.etoile}
\tag{$*$} \H^1(S,G)\isomto \H^1(\pi_1,\cal{G})\quoi,
\end{equation*}
where
% original p. 300
the right-hand side denotes the set of classes, up to
isomorphism, of principal homogeneous bundles under
$\cal{G}$ in the category of finite sets on which $\pi_1$
acts (there is moreover no need to specify: continuously),
a set which is made explicit in a well-known way as the
quotient set of the set $Z^1(\pi_1,G)$ of $1$-cocycles
$\varphi:\pi_1\to \cal{G}$ (satisfying $\varphi(1)=1$,
$\varphi(st)=\varphi(s)(s.\varphi(t))$) by the group $\cal{G}$, which
acts on it in a natural way).

\label{page-300}
An important case is that where $\pi_1$ acts trivially on~$\cal{G}$, \ie when $G$ is a completely
decomposed covering of~$S$, isomorphic to the sum of~$\cal{G}$
copies of~$S$; one then also writes $\H^1(S,\cal{G})$ instead
of~$\H^1(S,G)$, and this set is in one-to-one correspondence
\eqref{eq:XI.5.etoile} with $\H^1(\pi_1,\cal G)= \Hom(\pi_1,\cal{G})/$inner automorphisms of~$\cal{G}$. Note moreover that in this case, a principal
homogeneous bundle over~$S$ with group $G$ is nothing but a
\emph{principal covering} of~$S$ with group $\cal{G}$
(V~\Ref{V.2.7}), and the preceding one-to-one correspondence
is the one deduced from the correspondence between principal
coverings of~$S$ with group~$\cal{G}$, \emph{pointed} over~$a$, and the continuous homomorphisms from~$\pi_1(S,a)$ to $\cal{G}$
(\Ref{V},~end~of~\No \Ref{V.5}).

The interest of relating the theory of \'etale
coverings with that of principal bundles (already
implicit in \textsc{A\ptbl Weil}, G\'en\'eralisation des Fonctions
Ab\'eliennes, and made explicit by \textsc{S\ptbl Lang} in his
geometric class field theory, \cf Serre~\cite{XI.5}) comes from the fact
that every $S$-group that is finite and \'etale over~$S$ can be embedded
in an $S$-group~$H$, affine and smooth over~$S$, with
connected fibers, commutative when $G$ is, so that by the exact
sequence~\Ref{XI.4.5} (and possibly its noncommutative
variants), the ``discrete'' classification of principal
coverings with group $G$ can be studied with the help of the
``continuous'' classification of principal bundles with group~$H$,
and conversely, for that matter. For the idea of the
general construction of the immersion of~$G$ into $H$ (rather little
used in practice, it seems), refer to
\cite[VI~2.8]{XI.5}. We content ourselves with developing in the next \No
two important special cases, which are moreover classical. We shall
need there an auxiliary result:

\begin{proposition}
\label{XI.5.1}
Let
% original p. 301
$S$ be a prescheme, and $G$ an $S$-group isomorphic to
$\Gl(n)_S$ (for example $\GG_{m\,S}$) or to $\GG_{a\,S}$; then every
principal homogeneous bundle under $G$ is locally trivial.
\end{proposition}

Let us make precise that $\Gl(n)_S$ ($n$ an integer $\geq 0$) denotes the
$S$-group that represents the contravariant functor $T\mto
\Gl(n,\Gamma(T,\cal{O}_T))$ in the $S$-prescheme~$T$; in
particular $\GG_{m\,S}$ (``multiplicative group over~$S$'')
represents the contravariant functor $T\mto
\Gamma(T,\cal{O}_T^*)$, hence as a prescheme over~$S$ is
isomorphic to $\Spec {\cal{O}_S[t,t^{-1}]}$, where $t$ is an
indeterminate. Similarly $\GG_{a\,S}$ represents the
contravariant functor $T\mto \Gamma(T,\cal{O}_T)$; it is therefore
isomorphic as an $S$-prescheme to $\Spec(\cal{O}_S[t])$,
where $t$ is an indeterminate. Note that by d\'evissage,
\Ref{XI.5.1} gives back the local triviality result of
Rosenlicht, relating to the case where $G$ admits a ``composition
series'' whose consecutive factors are groups of the
type considered here. (For a finer study of the questions of
local triviality of principal homogeneous bundles,
\cf \cite{XI.7}~and~\cite{XI.3}).

The first assertion is proved by remarking that
$G(T)=\Aut(\cal{O}_T^n)$, and that the morphisms $S'\to S$ occurring
in~\Ref{XI.4.1} (\ie which are faithfully flat and
quasi-compact) are effective descent morphisms for the
fibered category of Modules locally isomorphic to
$\cal{O}_T^n$, \ie locally free of rank~$n$
(VIII~\Ref{VIII.1.12}). The second is proved in an
analogous way, by noting that in this case one has $G(T)=\Aut(\cal{E}_T)$,
where $\cal{E}_T$ is the trivial \emph{extension} of~$\cal{O}_T$ by
$\cal{O}_T$ (and where the automorphisms must of course
respect the extension structure), that the morphisms $S'\to S$
occurring in~\Ref{XI.4.1} are effective descent
morphisms for the fibered category of extensions of~$\cal{O}_T$ by $\cal{O}_T$ (as follows easily from
VIII~\Ref{VIII.1.1}), and that such extensions are automatically
locally trivial.

\begin{remark}
\label{XI.5.2}
Note that the same type of proof applies to the
symplectic group $\Symp(2n)_S$, taking into account that an alternating
form on a module locally isomorphic to
$\cal{O}_S^{2n}$ which is ``nondegenerate'',
\ie defines an isomorphism of this Module onto its dual, is
locally isomorphic to the standard form. The result for the
orthogonal group is on the other hand false, already if $S$ is the
spectrum of a field, because there can be
% original p. 302
quadratic forms over a field that are not isomorphic to the
standard form. Moreover, one essentially shows in \cite{XI.3} that the
groups $\Gl$, $\Symp$, $\GG_a$ and those that decompose by d\'evissage into such
groups are, more or less, the only ones for which one
has a local triviality result of the type considered
here.
\end{remark}

\begin{corollary}
\label{XI.5.3}
One has canonical bijections
$$
\H^1(S,\GL(n)_S)\isomfrom \H^1(S,\GL(n,\cal{O}_S))\quoi,
$$
in particular
$$
\H^1(S,\GG_{m\,S})\isomfrom \H^1(S,\cal{O}_S^*)\quoi,
$$
and
$$
\H^1(S,\GG_{a\,S})\isomfrom \H^1(S,\cal{O}_S)\quoi,
$$
where the right-hand sides denote cohomologies of
the topological space $S$ with coefficients in ordinary
sheaves.
\end{corollary}

In particular, $\H^1(S,\GL(n)_S)$ is identified with the set of
classes, up to isomorphism, of locally free Modules of
rank $n$ over~$S$, and $\H^1(S,\GG_{a\,S})$ is identified with the set
of classes of extensions of the Module $\cal{O}_S$ by itself.
